%! Exponential Function Context and Data Modeling — Unit 4, Lesson 4.6 — Exit Ticket
\documentclass[10pt]{article}
\usepackage{precalculus-article}
\usepackage{precalculus-boxes}
\newcommand{\TallMath}[1]{$\displaystyle #1\rule[-1.4em]{0pt}{3.2em}$}

\begin{document}

\pageheader{Unit 4, Lesson 4.6}{Exit Ticket}
\namedateperiod

\vspace{0.14in}

\begin{enumerate}[itemsep=16pt]

\item A population grows at 6\% per year and starts at 5{,}000.

      Write the model $P(t)$ in $ab^t$ form: \quad $P(t) = $ \blank{6cm}

      \vspace{0.06in}
      Find $P(10)$. Round to the nearest whole number.

      \vspace{0.06in}
      $P(10) = $ \blank{7cm}

\item An exponential model passes through $(1,\,6)$ and $(3,\,54)$. Find $a$ and $b$.

      \vspace{0.06in}
      Show your work:

      \vspace{0.5in}

      $a = $ \blank{3cm} \quad $b = $ \blank{3cm}

\item For the model in \#1, what does the base $b = 1.06$ tell you about the population?
      Answer in a complete sentence.

      \vspace{0.06in}
      \writelines{2}

\end{enumerate}

\end{document}
